\documentclass[11pt,a4paper]{article}

\usepackage[T1]{fontenc}
\usepackage[utf8]{inputenc}
\usepackage{lmodern}
\usepackage[margin=2.5cm]{geometry}
\usepackage{amsmath,amssymb,amsthm}
\usepackage{booktabs}
\usepackage{array}
\usepackage{enumitem}
\usepackage{xcolor}
\usepackage[colorlinks=true,linkcolor=blue!50!black,citecolor=blue!50!black,urlcolor=blue!50!black]{hyperref}

\usepackage{microtype}
\setlength{\emergencystretch}{2em}
\newcommand{\code}[1]{\texttt{#1}}
\newcommand{\kt}{\mathbf{k}_t}
\newcommand{\kn}{\mathbf{k}_n}
\newcommand{\km}{\mathbf{k}_m}
\newcommand{\bfrho}{\boldsymbol{\rho}}
\newcommand{\Ghat}{\widehat{G}}
\newcommand{\dd}{\mathrm{d}}
\newcommand{\ii}{\mathrm{i}}
\newcommand{\sinc}{\operatorname{sinc}}
\newcommand{\abs}[1]{\left\lvert #1 \right\rvert}
\newcommand{\norm}[1]{\left\lVert #1 \right\rVert}

\newtheorem{proposition}{Proposition}
\theoremstyle{remark}
\newtheorem{remark}{Remark}

\title{Step 3 of the precomputed Green's operator:\\
z skeletons for the cross-layer and biaxial-scattered blocks}
\author{YellowThinFilms --- branch \code{z-skeletons}}
\date{16 October 2026}

\begin{document}
\maketitle

\begin{abstract}
The precomputed operator of \code{mod\_GreenOperator.jl} maps dipole moments on the
voxel planes of a source layer to the field on the planes of every observation layer.
Same-layer blocks in isotropic/uniaxial layers are Toeplitz + Hankel in the plane
indices and are applied by a z-FFT. Cross-layer blocks, and the scattered part of a
biaxial same-layer block, have no such structure; until now they were stored for every
(source plane, observation plane) pair. This document derives what step 3 replaces them
with, starting from the $\kt$-integral the code evaluates. In a homogeneous layer the
z-dependence of every plane-wave component is an exponential $e^{\pm\ii Q z}$ fixed by
the layer alone, so the field on all planes of a layer is interpolated from a few
\emph{skeleton} planes with \emph{one} matrix, the same for every quadrature node.
Because the matrix does not depend on the node, it commutes with the $\kt$-quadrature and
the lateral FFT, and the operator needs the integral only on the skeleton pairs. We give
the construction (an interpolative decomposition checked on every node), its error
bound, the measured sizes, the apply algorithm and the validation.
\end{abstract}

\tableofcontents

%======================================================================
\section{The integral being evaluated}
\label{sec:integral}
%======================================================================

\subsection{Master identity}

A source with moment $\mathbf{p}$ at depth $z$ in the source layer $s$, observed at depth
$z'$ in layer $o$ and lateral offset $\bfrho=(\Delta x,\Delta y)$, produces the field
\begin{equation}
  E_c(\bfrho; z', z) \;=\; \sum_{q=1}^{3} G_{cq}(\bfrho; z', z)\, p_q,
  \qquad
  G_{cq}(\bfrho; z', z) \;=\; \iint \frac{\dd^2 \kt}{(2\pi)^2}\;
     S(\kt)\, e^{\ii \kt\cdot\bfrho}\; \widetilde{G}_{cq}(\kt; z', z),
  \label{eq:master}
\end{equation}
with $\kt=(k_\rho\cos\varphi, k_\rho\sin\varphi)$ conserved through the stack. $S$ is the
lateral footprint factor of the source: $S\equiv 1$ for a point, and for a rectangular
box of size $\Delta x_b\times\Delta y_b$
\begin{equation}
  S(\kt) = \sinc\!\left(\tfrac{k_x\Delta x_b}{2}\right)\sinc\!\left(\tfrac{k_y\Delta y_b}{2}\right),
  \qquad \sinc u = \frac{\sin u}{u},
\end{equation}
the exact Fourier transform of the footprint (\code{box\_lateral\_factor}).

\subsection{Separation in $(z', z)$ for one plane wave}

For one $\kt$, each layer carries four Berreman eigenmodes: two forward ($f$) and two
backward ($b$), with normal wavenumbers $Q_{f,1}, Q_{f,2}, Q_{b,1}, Q_{b,2}$
($\operatorname{Im} Q_f \ge 0$, $\operatorname{Im} Q_b \le 0$). Write $Q^{(s)}$ for the
source layer's wavenumbers and $K^{(o)}$ for the observation layer's. The response
separates exactly (\code{mod\_CavityTransfer.jl}, \code{mod\_ClosedFormField.jl}):
\begin{align}
  \widetilde{G}_{cq}(\kt; z', z)
  &= \widetilde{G}^{\,\mathrm{dir}}_{cq}(\kt; z', z)\,\delta_{os}
   + \sum_{\beta\in\{f,b\}}\sum_{m=1}^{2}
     e^{(o)}_{c;\beta m}(\kt)\;\Phi_{\beta m}(\kt; z')\; A_{\beta m, q}(\kt; z),
  \label{eq:separation}\\
  A_{\beta m, q}(\kt; z)
  &= \sum_{\gamma\in\{f,b\}}\sum_{\nu=1}^{2}
     T^{(o)}_{\beta m;\gamma\nu}(\kt)\; s_{\gamma\nu, q}(\kt)\; D_{\gamma\nu}(\kt)\; \Psi_{\gamma\nu}(\kt; z).
  \label{eq:amplitude}
\end{align}
The pieces are:
\begin{itemize}[leftmargin=*]
\item \textbf{Source side.} With $z_{\mathrm{pos}}$ the cavity-frame depth ($z_{\mathrm{pos}} = d_s - z$ in layer 1,
  $z_{\mathrm{pos}}=z$ otherwise) and $d_s$ the source layer's thickness,
  \begin{equation}
    \Psi_{f\nu}(z) = e^{\ii Q^{(s)}_{f,\nu}(d_s - z_{\mathrm{pos}})},
    \qquad
    \Psi_{b\nu}(z) = e^{-\ii Q^{(s)}_{b,\nu}\, z_{\mathrm{pos}}}.
    \label{eq:psi}
  \end{equation}
  $s_{\gamma\nu,q}$ is the point seed of moment $\hat{\mathbf{e}}_q$ (\code{Source\_projection\_Berreman\_xyz}),
  and $D_{\gamma\nu} = \sinc(Q^{(s)}_{\gamma,\nu}\Delta z_b/2)$ the axial box factor ($D=1$ for a point).
  The code forms $D\Psi$ together as a face difference, but mathematically it is the
  product of a per-node constant and the exponential~\eqref{eq:psi} at the box centre.
\item \textbf{Stack.} $T^{(o)}(\kt)$ is a $4\times4$ matrix per node and observation layer: the
  cavity resummation and the walk to layer $o$. It depends on neither $z$ nor $z'$.
\item \textbf{Observation side.} $e^{(o)}_{c;\beta m}$ is component $c$ of the mode field of mode $(\beta,m)$, and
  \begin{equation}
    \Phi_{f m}(z') = e^{\ii K^{(o)}_{f,m} z'},
    \qquad
    \Phi_{b m}(z') = e^{\ii K^{(o)}_{b,m}(z' - d_o)}.
    \label{eq:phi}
  \end{equation}
\item \textbf{Direct term.} $\widetilde{G}^{\,\mathrm{dir}}$ is the source's own radiation in its layer: forward
  branch for $z' > z_{\mathrm{pos}}$, backward branch below, and face ramps plus the local $E_z$
  inside a box's own slab. It is the only non-separable piece: it contains a Heaviside
  step in $z'-z_{\mathrm{pos}}$.
\end{itemize}

\begin{proposition}[Bounded exponentials]
\label{prop:bounded}
All of $\abs{\Psi_{\gamma\nu}(z)}$ and $\abs{\Phi_{\beta m}(z')}$ are $\le 1$ for planes inside their layers.
\end{proposition}
\begin{proof}
Each exponent is $\ii Q\ell$ with $\ell$ the distance from the plane to the interface
the wave is referenced to, and the sign of $\operatorname{Im} Q$ makes the wave decay
away from that interface. Forward: $\operatorname{Im} Q_f\ge0$ and $\ell=d_s-z_{\mathrm{pos}}\ge0$, or
$\ell=z'\ge0$. Backward: $\operatorname{Im} Q_b\le0$ and $\ell=-z_{\mathrm{pos}}\le0$, or
$\ell=z'-d_o\le 0$. In each case $\operatorname{Re}(\ii Q\ell)\le0$.
\end{proof}

\subsection{Quadrature}

The code evaluates \eqref{eq:master} on a $(k_\rho,\varphi)$ rule: nodes $\kn$ with weights
$w^\rho_n w^\varphi_n$, adaptive in $k_\rho$ (\code{box\_k$\rho$\_grid}, or one $k_\rho$ rule per
$\varphi$ group, \code{box\_k$\rho$\_grid\_grouped}). The node weight used in \code{closed\_form\_G} is
\begin{equation}
  W_n = \frac{\ii\,\omega\, k_{\rho,n}\, w^\rho_n\, w^\varphi_n}{8\pi^3}\; S(\kn),
  \label{eq:weights}
\end{equation}
that is, the polar Jacobian $k_\rho$ and the $1/(2\pi)^2$ of \eqref{eq:master}, times the
code's prefactor $\ii\omega/(2\pi)$. Let $g_n(z',z)\in\mathbb{C}^{3\times3}$ be the
separable part of $\widetilde G(\kn; z', z)$, i.e.\ the double sum in \eqref{eq:separation}.
The discrete field is then
\begin{equation}
  G(\bfrho; z', z) \;\approx\; \sum_n W_n\, e^{\ii\kn\cdot\bfrho}\, g_n(z', z)
  \;+\; (\text{direct term, same layer}).
  \label{eq:discrete}
\end{equation}
The lateral sum is done exactly or by a type-1 NUFFT (accuracy \code{nufft\_eps}).
Both are linear in the per-node strengths $W_n g_n$, which is all that is used below.

\paragraph{Rank $\le 4$ per node.}
Collect the four source-side exponentials in $\psi_n(z) = (\Psi_{f1},\Psi_{f2},\Psi_{b1},\Psi_{b2})^{\mathsf T}$
and the four observation-side ones in $\phi_n(z')$. Then \eqref{eq:separation}--\eqref{eq:amplitude} read
\begin{equation}
  g_{n,cq}(z',z) = \phi_n(z')^{\mathsf T}\, C_{n,cq}\, \psi_n(z),
  \qquad
  \bigl(C_{n,cq}\bigr)_{\beta m,\gamma\nu} = e^{(o)}_{c;\beta m}\,T^{(o)}_{\beta m;\gamma\nu}\,s_{\gamma\nu,q}\,D_{\gamma\nu},
  \label{eq:outer}
\end{equation}
with $C_{n,cq}\in\mathbb{C}^{4\times4}$. Every source--observation dependence of one node passes through four
edge amplitudes, so on any set of planes $g_n$ has rank $\le 4$. It has rank $\le2$ into a
semi-infinite layer, where only outgoing waves exist (\S\ref{sec:onesided}).

%======================================================================
\section{From the integral to the precomputed operator}
\label{sec:operator}
%======================================================================

\subsection{Lateral FFT on the padded offsets}

The sources sit on an $N_x\times N_y$ pixel grid with spacings $\Delta x,\Delta y$, on planes
$z_1<\dots<z_{N_s}$ of the source layer. The observation planes of layer $o$ are
$z'_1<\dots<z'_{N_o}$. $G$ is sampled on the $P_x\times P_y=2N_x\times 2N_y$ padded offsets
$\bfrho_{ab}=(a\Delta x, b\Delta y)$, $a=-N_x,\dots,N_x-1$ (\code{padded\_offsets}). Its
transform is stored at the padded frequencies
$\km = 2\pi\,(m_x/(P_x\Delta x),\, m_y/(P_y\Delta y))$:
\begin{equation}
  \Ghat_{ji}(\km) = \frac{1}{P_xP_y}\sum_{a,b} G(\bfrho_{ab}; z'_j, z_i)\, e^{-\ii\km\cdot\bfrho_{ab}}
  \qquad(\code{\_ghat\_planes}).
  \label{eq:ghat}
\end{equation}
With $\widehat{p}_i = \mathrm{FFT}(p_i)$ on the zero-padded grid, the field is
\begin{equation}
  \widehat{E}_j(\km) = \sum_{i=1}^{N_s} \Ghat_{ji}(\km)\,\widehat{p}_i(\km),
  \qquad E_j = \mathrm{BFFT}(\widehat{E}_j)\big|_{N_x\times N_y},
  \label{eq:apply}
\end{equation}
which is the exact linear convolution $\sum_i\sum_{\mathbf{x}'} G(\mathbf{x}-\mathbf{x}'; z'_j,z_i)\,p_i(\mathbf{x}')$:
every source--observation offset is one of the $\bfrho_{ab}$, and $\mathrm{BFFT}\circ\mathrm{FFT}=P_xP_y$.

\subsection{What a pixel frequency sees}

Insert the node sum \eqref{eq:discrete} into \eqref{eq:ghat}:
\begin{equation}
  \Ghat_{ji}(\km) = \sum_n W_n\, \mathcal{K}(\kn - \km)\, g_n(z'_j, z_i),
  \qquad
  \mathcal{K}(\mathbf{q}) = \frac{1}{P_xP_y}\sum_{a,b} e^{\ii\mathbf{q}\cdot\bfrho_{ab}},
  \label{eq:kernel}
\end{equation}
plus the direct term. $\mathcal{K}$ is a two-dimensional Dirichlet kernel. Its main lobe has width
$2\pi/(P_x\Delta x)=2\pi/(4L)$ for a window of half-width $L$, so $\Ghat(\km)$ is a
weighted sum over a \emph{bundle} of nodes around $\km$, plus their aliases $\km+2\pi\mathbf{n}/\Delta$.
This truncation to the window is what makes \eqref{eq:apply} exact on the window, so it is not an error to be removed.

\subsection{Block structure before step 3}

Per (source layer, observation layer) the operator stores one block (\code{ZBlock}):
\begin{itemize}[leftmargin=*]
\item \textbf{Same layer, isotropic or uniaxial (azimuthally symmetric stack).} $g_n(z',z)$ depends on $z'-z_{\mathrm{pos}}$ (direct + matched-branch
  reflections) and on $z'+z_{\mathrm{pos}}$ (cross-branch reflections) only. On uniform planes the block is
  Toeplitz + Hankel, so $N_s+N_o-1$ planes each, applied by a z-FFT (\code{:structured}).
\item \textbf{Same layer, biaxial.} The interfaces mix the two modes, which have different $Q$, so only the
  direct term stays Toeplitz. The scattered part was stored for all $N_s N_o$ pairs (\code{:toeplitz\_dense}).
\item \textbf{Cross layer.} $Q^{(s)}\ne K^{(o)}$, so there is no translation structure: all $N_sN_o$ pairs (\code{:dense}).
\end{itemize}
A dense block holds $P_xP_y\cdot 9\cdot N_s N_o$ complex numbers and needs $N_s N_o$ evaluations of
\eqref{eq:discrete} (one \code{closed\_form\_G} row each). It was the memory and build bottleneck:
\(\approx\)100--140~s and 600~MB per block at the benchmark's size L.

\subsection{Why not four amplitudes per pixel frequency}
\label{sec:four}

By \eqref{eq:outer}, a node needs only $4(N_s+N_o)$ numbers: four edge amplitudes per source plane
and per observation plane. In \eqref{eq:kernel}, however, every node of the bundle has its own $\psi_n$ and $\phi_n$:
\begin{equation}
  \Ghat(\km) = \sum_n \underbrace{W_n\,\mathcal{K}(\kn-\km)}_{\text{scalar}}\;\phi_n\,C_n\,\psi_n^{\mathsf T},
\end{equation}
a sum of many rank-4 terms with different exponentials, so the result is not rank 4. Measured on the dense blocks
(\code{diagnostics/zrank\_study.jl}, 16 planes per layer, $\Delta=1/32$):
\begin{itemize}[leftmargin=*]
\item the best rank-4 approximation per $\km$ leaves $4\times10^{-2}$ of the cross block;
\item capping $k_\rho$ below the pixel band, so that nothing aliases, does not change this;
\item widening the window helps, but slowly.
\end{itemize}
The only exact version keeps the nodes in the apply: a type-2 NUFFT of every source plane from the pixels to the
$\sim10^6$ nodes, the contraction with \eqref{eq:outer}, and a type-1 NUFFT back. That is
$\sim 10^6/(P_xP_y)\approx 250$ times more points per plane than the lateral FFTs at every apply, so it is not used.

%======================================================================
\section{z skeletons}
\label{sec:skeleton}
%======================================================================

\subsection{The exponential family of a layer}

The idea is to compress what all nodes share, not each $\km$. Every node's dependence on the observation
planes of layer $o$ lies in the span of
\begin{equation}
  \mathcal{F}^{\mathrm{obs}}_o = \Bigl\{\, \bigl(e^{\ii K z'_j}\bigr)_{j=1}^{N_o} : K\in\{K^{(o)}_{f,m}(\kn)\}_{n,m} \Bigr\}
  \;\cup\;
  \Bigl\{\, \bigl(e^{\ii K (z'_j-d_o)}\bigr)_{j=1}^{N_o} : K\in\{K^{(o)}_{b,m}(\kn)\}_{n,m} \Bigr\},
  \label{eq:famobs}
\end{equation}
and its dependence on the source planes in the span of
\begin{equation}
  \mathcal{F}^{\mathrm{src}}_s = \Bigl\{\, \bigl(e^{\ii Q (d_s-z_{\mathrm{pos},i})}\bigr)_{i=1}^{N_s} : Q\in\{Q^{(s)}_{f,\nu}(\kn)\} \Bigr\}
  \;\cup\;
  \Bigl\{\, \bigl(e^{-\ii Q z_{\mathrm{pos},i}}\bigr)_{i=1}^{N_s} : Q\in\{Q^{(s)}_{b,\nu}(\kn)\} \Bigr\}.
  \label{eq:famsrc}
\end{equation}
Arrange each family as a matrix $F\in\mathbb{C}^{N\times M}$: one row per plane, one column per (node, mode, branch).
After removing duplicate $Q$ values, $M\sim 10^5$--$10^6$. The families depend on the layer's
dispersion and the planes only, \emph{not} on the stack. The stack (reflections, guided modes,
plasmons) lives entirely in the coefficients $C_n$ of \eqref{eq:outer}. The box factor $D$ and the seeds are
per-node constants and do not enter either.

\paragraph{Physical reading.} A plane wave entering a homogeneous layer keeps its $\kt$, so its z-profile there is set by the layer
alone. $F$ has two parts:
\begin{itemize}[leftmargin=*]
\item \textbf{Propagating part} ($Q$ real in $[0, n k_0]$, with $n$ the layer's index): a band-limited family. Its numerical
  dimension is the Nyquist count $\approx 2 n t/\lambda$ per direction, for a layer of thickness $t$, plus a
  term in $\log(1/\tau)$ (prolate spheroidal functions).
\item \textbf{Evanescent part} ($e^{-\kappa\ell}$ with $\kappa\in[\kappa_{\min},\kappa_{\max}]$): a family of decaying
  exponentials. These are nearly linearly dependent, and their numerical rank grows like
  $\log(\kappa_{\max}/\kappa_{\min})\cdot\log(1/\tau)$.
\end{itemize}
Neither part grows with the number of planes beyond logarithmically, so a fine plane grid is
highly redundant. The exception is a hyperbolic layer ($\varepsilon_{zz}<0<\varepsilon_t$), where
$Q\approx\sqrt{\varepsilon_t/\abs{\varepsilon_{zz}}}\,k_\rho$ is real up to the $k_\rho$ cap. High-$k_t$ waves
propagate there, the band limit exceeds the plane Nyquist frequency $\pi/\Delta$, and the family has full rank.

\subsection{One-sided families in semi-infinite layers}
\label{sec:onesided}

In the bottom layer (layer 1) no wave comes from $-\infty$, and in the top layer none from $+\infty$. So
\begin{itemize}[leftmargin=*]
\item observation layer 1 carries only the backward (downward) wave;
\item the last layer carries only the forward one;
\item a source in layer 1 feeds only through $\Psi_f$.
\end{itemize}
This was checked on a cache: the corresponding rows and columns of $T^{(o)}$ are exactly 0 for every node.
The families drop the other branch (\code{\_obs\_family}, \code{\_src\_family}), which roughly halves their rank.

\subsection{Interpolative decomposition}

An interpolative decomposition (ID) of $F$ on its rows is a set of $R$ \emph{skeleton planes}
$\mathcal{S}=\{\sigma_1<\dots<\sigma_R\}$ and a matrix $P\in\mathbb{C}^{N\times R}$ with $P_{\mathcal{S},:}=I$ such that
\begin{equation}
  \max_{j,\,\mu}\;\bigl|\,F_{j\mu} - \textstyle\sum_{r=1}^{R} P_{jr}\,F_{\sigma_r\mu}\bigr| \;\le\; \tau
  \qquad\text{for \emph{every} column } \mu.
  \label{eq:id}
\end{equation}
By Proposition~\ref{prop:bounded}, every entry of $F$ has modulus $\le1$, so $\tau$ is an absolute bound
\emph{per node term}. In words: the value of every plane wave of the layer on every plane is reproduced, to
$\tau$, from its values on the skeleton planes, with the same weights $P_{jr}$ for all waves.

\subsection{Substitution into the integral}

Let $(\mathcal S_s, P^s)$ and $(\mathcal S_o, P^o)$ be the IDs of \eqref{eq:famsrc} and \eqref{eq:famobs}. Applied to
\eqref{eq:outer}, they give for every node, component pair $(c,q)$ and plane pair $(j,i)$
\begin{equation}
  g_n(z'_j, z_i) \;=\; \sum_{r=1}^{R_o}\sum_{r'=1}^{R_s} P^o_{jr}\,P^s_{ir'}\; g_n(z'_{\sigma^o_r}, z_{\sigma^s_{r'}})
  \;+\; \varepsilon_{n,ji}.
  \label{eq:node_interp}
\end{equation}
$P^o$ and $P^s$ do not depend on $n$. Multiplying by $W_n e^{\ii\kn\cdot\bfrho}$ and summing over $n$, then
transforming by \eqref{eq:ghat}, therefore gives the same relation for the field and for its FFT:
\begin{equation}
  \boxed{\;\Ghat_{ji}(\km) \;=\; \sum_{r,r'} P^o_{jr}\,P^s_{ir'}\;\Ghat_{\sigma^o_r\,\sigma^s_{r'}}(\km) \;+\; \mathcal{E}_{ji}(\km)\;}
  \label{eq:main}
\end{equation}
The interpolation commutes with the $\kt$-quadrature, the NUFFT (linear in the strengths), the window kernel
$\mathcal K$ and the lateral FFT. So the block is fully determined by the integral \eqref{eq:master} evaluated
on the $R_sR_o$ skeleton pairs only. This is the step that fails per $\km$ (\S\ref{sec:four}) and works per node.
The same holds for the grouped build: the blocks of the $\varphi$ groups are summed, and \eqref{eq:main} holds
group by group \emph{provided} all groups use the same $(\mathcal S, P)$.

\begin{proposition}[Error bound]
\label{prop:error}
Write $\norm{P}_\infty=\max_j\sum_r\abs{P_{jr}}$ and $\norm{C}_1=\sum\abs{C_{\beta m,\gamma\nu}}$. Under \eqref{eq:id} for both sides,
\begin{equation}
  \abs{\varepsilon_{n,ji}} \le \tau\,\bigl(1+\norm{P^o}_\infty\bigr)\,\norm{C_n}_1,
  \qquad
  \abs{\mathcal{E}_{ji}(\km)} \le \tau\,\bigl(1+\norm{P^o}_\infty\bigr)\sum_n \abs{W_n\,\mathcal{K}(\kn-\km)}\,\norm{C_n}_1 .
\end{equation}
\end{proposition}
\begin{proof}
Write $\phi=P^o\phi_{\mathcal S}+\delta^o$ and $\psi=P^s\psi_{\mathcal S}+\delta^s$, with $\abs{\delta}\le\tau$ entrywise. Then
$\phi^{\mathsf T}C\psi - (P^o\phi_{\mathcal S})^{\mathsf T}C(P^s\psi_{\mathcal S})
= (P^o\phi_{\mathcal S})^{\mathsf T}C\,\delta^s + (\delta^o)^{\mathsf T}C\,\psi$.
Use $\abs{P^o\phi_{\mathcal S}}\le\norm{P^o}_\infty$ and $\abs{\psi}\le1$, then sum over nodes.
\end{proof}

\begin{remark}[What the bound does and does not say]
The bound is a priori and not sharp: it ignores cancellation between nodes and is relative to
$\sum_n\abs{W_n C_n}$, not to $\abs{G}$.
For a box, $D$ is folded into $C_n$ and grows like $e^{\abs{\operatorname{Im}Q}\Delta z_b/2}$ for deep-evanescent nodes,
so the bound becomes weak exactly there. In practice the ID places skeleton planes next to both interfaces.
For a glass layer at 64 planes and $\tau=10^{-6}$ they are planes $1$--$6$ and $59$--$64$, with sparser ones towards the middle:
the evanescent waves vary fastest near the interfaces. The deep-evanescent columns are therefore reproduced
\emph{exactly} on the planes where they are not negligible. Measured on that family: $\max_\mu \abs{\delta_\mu}\cdot\abs{D_\mu} = 6.6\times10^{-7}$,
against $\max_\mu\abs{\delta_\mu} = 5.2\times10^{-7}$. The box factor does not degrade the interpolation.
The field errors measured against the dense blocks are $0.1$--$0.7\,\tau$ relative (\S\ref{sec:validation}).
\end{remark}

%======================================================================
\section{Construction}
\label{sec:construction}
%======================================================================

\subsection{Algorithm (\code{\_zskeleton})}

\begin{enumerate}[leftmargin=*]
\item \textbf{Exact compression of the columns.} Form $F^{\mathsf T}$ (nodes $\times$ planes) in chunks of 8192 rows and reduce it to its
  $N\times N$ triangular factor $R$ by a threaded TSQR ($F^{\mathsf T} = Q_0R$). Since $Q_0$ has orthonormal columns,
  a column ID of $R$ is a column ID of $F^{\mathsf T}$, with the same error in every norm. No column is
  sampled, so the extreme evanescent nodes cannot be missed.
\item \textbf{Pivot order.} Column-pivoted QR of $R$: $R\Pi = \widetilde Q\widetilde R$. For a rank $k$, the skeleton is the
  first $k$ pivots, and $P$ follows from $\widetilde R_{11}^{-1}\widetilde R_{12}$, with the identity on the skeleton rows.
\item \textbf{Rank chosen on every node.} For a candidate $k$, evaluate the left side of \eqref{eq:id} over all columns,
  streamed again in chunks. Bisect for the smallest $k$ that satisfies \eqref{eq:id} ($k=N$ is exact).
\item \textbf{Fallback.} If $k\ge0.8\,N$, keep the side whole ($\mathcal S=\{1..N\}$, $P=I$, stored as \code{nothing}).
  If both sides of a block are whole, the block stays the plain \code{:dense} / \code{:toeplitz\_dense} one.
\end{enumerate}
\textbf{Cost.} About 2~s for 750k columns at 64 planes on 8 threads. BLAS is pinned to one thread inside: the
chunk loops are threaded, and nested BLAS threads were $\approx10\times$ slower.
\textbf{Interpolation norm.} $\norm{P}_2 = 1.8$ (16 planes) and $3.2$--$3.7$ (64 planes), with $\norm{P}_\infty\le 4.7$. The
interpolation does not amplify errors.

\subsection{Where the $Q$'s come from}

\begin{itemize}[leftmargin=*]
\item \textbf{Single cache} (\code{build\_layer\_operator}). The cache's own \code{Qf\_src}, \code{Qb\_src} and
  \code{basis[o].KZf}, \code{KZb} are used: these are exactly the $Q$'s in \eqref{eq:psi}--\eqref{eq:phi}.
\item \textbf{Grouped build} (\code{build\_layer\_operator\_grouped}). The blocks are summed over $\varphi$ groups and $\varphi$ chunks, so
  all must share one skeleton (\S\ref{sec:skeleton}). \code{z\_skeletons\_grouped} builds it before the first cache, from
  each layer's own mode solve (\code{\_modes\_layer}, a $4\times4$ eigenproblem, not a stack solve) at every $(k_\rho,\varphi)$
  node of every group. On an azimuthally symmetric stack only $\varphi=0$ is used, since $Q$ depends on $k_\rho$ alone.
  The set of $Q$'s is the same as in the caches: the stack solver only reorders modes.
\end{itemize}
The source skeleton is built in the operator's internal plane order (increasing $z_{\mathrm{pos}}$), so the
reversed planes of layer 1 need no special handling.

\subsection{Blocks}

\begin{itemize}[leftmargin=*]
\item \code{:skeleton} (cross layer) stores $\Ghat_{\sigma^o_r\sigma^s_{r'}}(\km)$ as \code{(Px, Py, 3, 3, R\_s, R\_o)}, plus $P^s$ and/or $P^o$.
  The table is evaluated with the existing \code{\_dense\_planes} on the skeleton planes: no new kernel code.
\item \code{:toeplitz\_skeleton} (biaxial same layer): the Toeplitz direct term, exactly as before, plus the scattered part as a
  skeleton block. The direct term is not compressed: its Heaviside branch in $z'-z_{\mathrm{pos}}$ is not of the form
  \eqref{eq:outer}, and it is already Toeplitz.
\item Isotropic/uniaxial same-layer blocks are unchanged (\code{:structured}).
\item \code{structure=:dense} disables the skeletons; it is the reference.
\end{itemize}

%======================================================================
\section{Apply}
\label{sec:apply}
%======================================================================

Equation \eqref{eq:apply} with \eqref{eq:main} becomes, per $\km$ and component:
\begin{align}
  \widehat{p}^{\,\mathrm{sk}}_{r'}(\km) &= \sum_{i=1}^{N_s} P^s_{ir'}\,\widehat p_i(\km), \label{eq:fold}\\
  \widehat{E}^{\,\mathrm{sk}}_{r}(\km) &= \sum_{r'=1}^{R_s} \Ghat_{\sigma^o_r\sigma^s_{r'}}(\km)\,\widehat{p}^{\,\mathrm{sk}}_{r'}(\km), \label{eq:pairs}\\
  \widehat{E}_j(\km) &\mathrel{+}= \sum_{r=1}^{R_o} P^o_{jr}\,\widehat{E}^{\,\mathrm{sk}}_r(\km). \label{eq:expand}
\end{align}
These are: \eqref{eq:fold} the fold of the source planes (\code{\_zfold}, one GEMM per component over all $\km$);
\eqref{eq:pairs} the pair loop on the skeleton planes (\code{\_dense\_add!}); \eqref{eq:expand} the interpolation to
all observation planes (\code{\_zexpand\_add!}, one GEMM per component).
followed by the usual 2D inverse FFT. The fold depends only on the source skeleton, so it is computed once and
shared by every block of the operator. Note that $\sum_i P^s_{ir'}\widehat p_i$ is the
\emph{transpose} of the interpolation: \eqref{eq:node_interp} interpolates $G$ in its source argument, so the sum over sources
$\sum_i G(\cdot,z_i)\,p_i$ becomes $\sum_{r'} G(\cdot,z_{\sigma^s_{r'}})\,(P^{s\mathsf T}p)_{r'}$.

\textbf{Costs per $\km$}, compared with a dense block:
\begin{center}
\begin{tabular}{lcc}
\toprule
 & dense & skeleton \\
\midrule
memory (complex numbers) & $9N_sN_o$ & $9R_sR_o + (N_sR_s + N_oR_o)/(P_xP_y)$ \\
build (\code{closed\_form\_G} rows) & $N_sN_o$ & $R_sR_o$ \\
apply (multiply--adds) & $9N_sN_o$ & $3N_sR_s + 9R_sR_o + 3R_oN_o$ \\
\bottomrule
\end{tabular}
\end{center}
The grid ($k_\rho$ rule) and the cache (stack solve) are unchanged by step 3. The precompute therefore drops by
less than the block ratio when those dominate.

%======================================================================
\section{Measured skeleton sizes}
\label{sec:ranks}
%======================================================================

The table below gives the rank $k$ chosen by \eqref{eq:id} at $\tau = 10^{-4}/10^{-6}/10^{-9}$, without the 80\% rule
(\code{diagnostics/zskeleton\_rank.jl}). The families use 6k $k_\rho\in[0,3000]$, $\times\,32$ $\varphi$ for the biaxial layers, both modes.
\begin{center}
\footnotesize
\setlength{\tabcolsep}{4pt}
\begin{tabular}{lccccccc}
\toprule
 & \multicolumn{5}{c}{$t=0.5\lambda$, $N$ planes} & \multicolumn{2}{c}{$t=2\lambda$} \\
\cmidrule(lr){2-6}\cmidrule(lr){7-8}
layer & 8 & 16 & 32 & 64 & 128 & 64 & 128 \\
\midrule
vacuum cladding (outgoing) & 8/8/8 & 9/12/15 & 11/14/18 & 12/16/21 & 13/17/24 & 14/18/23 & 16/20/26 \\
glass, obs = src & 8/8/8 & 12/15/16 & 15/19/24 & 18/23/30 & 20/27/36 & 21/25/33 & 23/30/38 \\
biaxial $\varepsilon_{zz}=+2$ & 8/8/8 & 12/15/16 & 16/19/24 & 18/24/30 & 21/27/37 & 24/27/34 & 26/32/40 \\
hyperbolic $\varepsilon_{zz}=-3.2$ & full & full & full & full & full & full & full \\
\bottomrule
\end{tabular}
\end{center}
\begin{itemize}[leftmargin=*]
\item The rank grows logarithmically in $N$, as expected from the evanescent part.
\item A 2$\lambda$ layer costs only 2--3 planes more than a 0.5$\lambda$ one at the same $N$: the propagating (Nyquist) part is
  small next to the evanescent one.
\item The biaxial family matches the glass one to $\pm1$ plane. The hyperbolic one is full rank, as predicted.
\item The earlier estimate (\code{diagnostics/zfamily\_rank.jl}) counted singular values relative to $\sigma_1$ of a sampled family.
  That under-counted by 1--4 planes, because $\sigma_1$ grows with the number of samples. Criterion \eqref{eq:id} does not
  depend on sampling.
\end{itemize}
Expected size of a vacuum $\leftrightarrow$ glass cross block at $\tau=10^{-6}$, $R_oR_s/(N_oN_s)$:
\begin{itemize}[leftmargin=*]
\item 0.75 at 16 planes (the glass side is kept whole by the 80\% rule);
\item 0.26 at 32;
\item 0.09 at 64 (11$\times$);
\item 0.028 at 128 (36$\times$).
\end{itemize}
The benchmark's present size (16 planes per layer) gains little. The gain is for fine $\Delta$ or many planes.

%======================================================================
\section{Validation}
\label{sec:validation}
%======================================================================

\code{Gz\_ClosedForm\_Operator\_Validation.jl}, all sections PASS:
\begin{itemize}[leftmargin=*]
\item \textbf{Exactness.} The tests [o] (operator vs the per-depth \code{closed\_form\_field} route), [g] (grouped $\varphi$ grid)
  and [m] (cost) run with \code{skeleton=false} and are unchanged: $\le 1.3\times10^{-15}$; a shared grid split into groups agrees to $6\times10^{-16}$.
\item \textbf{New [s].} Skeleton vs dense blocks for random moments on every plane, every observation layer,
  $\Delta=1/32$, 0.75$\lambda$ claddings (24 planes), 0.5$\lambda$ middle layer (16 planes), $\tau=10^{-6}$:
\end{itemize}
\begin{center}
\small
\begin{tabular}{>{\raggedright\arraybackslash}p{4.6cm}>{\raggedright\arraybackslash}p{6.8cm}c}
\toprule
case & blocks ($R_s\times R_o$) & rel.\ error \\
\midrule
iso box 2D, source layer 2 & $2\to1$: $16\times14$, $2\to2$: structured, $2\to3$: $16\times14$ & $7.6\times10^{-8}$ \\
iso box 2D, source layer 1 (reversed planes) & $1\to2$: $14\times16$, $1\to3$: $14\times14$ & $7.7\times10^{-8}$ \\
biaxial box 2D, layer 2 & $2\to1,3$: $16\times13$; $2\to2$ both sides whole & $4.4\times10^{-7}$ \\
hyperbolic box 2D, layer 2 & source side whole on every block (fallback) & $3.4\times10^{-7}$ \\
iso point, 1D cache, layer 2 & $2\to1,3$: $16\times13$ & $6.8\times10^{-7}$ \\
grouped biaxial box, 3 groups & one shared skeleton & $4.4\times10^{-7}$ \\
\bottomrule
\end{tabular}
\end{center}
The error is $\norm{E-E_{\mathrm{dense}}}/\norm{E_{\mathrm{dense}}}$, maximised over observation layers.
A smaller case (16 planes, $L=0.25$) gives field errors of $\approx 0.1\,\tau$ at $\tau=10^{-4}$ and $10^{-6}$.

%======================================================================
\section{Open points}
%======================================================================
\begin{itemize}[leftmargin=*]
\item \textbf{Benchmark.} \code{Bench\_GreenOperator.jl} now has a 2$\lambda$ source-layer stack (\code{thick}, 64 source planes at size L) and
  \code{BENCH\_SKELETON=0/1}. It has not been run yet.
\item \textbf{Default $\tau=10^{-6}$.} This is 100$\times$ below the grouped grid's \code{rtol} $=10^{-4}$. Whether it should follow \code{rtol}
  automatically is to be decided.
\item \textbf{Hyperbolic layers.} Their own side cannot be compressed by any node-independent z basis (the band limit exceeds $\pi/\Delta$).
  Only the other side of their cross blocks is compressed.
\end{itemize}

\appendix
\section{Code map}
\begin{center}
\small
\begin{tabular}{>{\raggedright\arraybackslash}p{6.2cm}>{\raggedright\arraybackslash}p{8.3cm}}
\toprule
symbol & code (\code{mod\_GreenOperator.jl} unless noted) \\
\midrule
$W_n$, \eqref{eq:weights}; $g_n$, \eqref{eq:discrete} & \code{closed\_form\_G}, \code{\_strengths\_row!} (\code{mod\_ClosedFormField.jl}) \\
$\Psi$, \eqref{eq:psi}; $D\Psi$ & \code{\_seed\_phases}, \code{\_boxphase} (\code{mod\_ClosedFormField.jl}) \\
$\Phi$, \eqref{eq:phi} & \code{\_obs\_phases}, \code{\_strengths\_row!} \\
$T^{(o)}$, $s$, $Q$, $K$ & \code{CavityTransferCache2D} / \code{1D} (\code{mod\_CavityTransfer.jl}) \\
$\Ghat(\km)$, \eqref{eq:ghat} & \code{\_ghat\_planes}, \code{\_dense\_planes} \\
families \eqref{eq:famobs}--\eqref{eq:famsrc} & \code{\_obs\_family}, \code{\_src\_family}, \code{\_family\_rows}, \code{\_qlist} \\
ID \eqref{eq:id} & \code{\_zskeleton} (TSQR via \code{\_fold\_chunks}, pivoted QR, bisection), \code{ZSkeleton} \\
skeletons of one operator & \code{z\_skeletons}, \code{\_cache\_skeletons}, \code{z\_skeletons\_grouped}, \code{\_layer\_Q} \\
blocks & \code{ZBlock} (fields \code{Ps}, \code{Po}), \code{\_planes\_maybe\_skeleton}, \code{\_dense\_block}, \code{\_toeplitz\_dense\_block} \\
apply & \code{apply\_operator}, \code{\_zfold}, \code{\_dense\_add!}, \code{\_zexpand\_add!} \\
keywords & \code{build\_layer\_operator(...; skeleton=true, skeleton\_tol=1e-6, skeletons=nothing)} \\
\bottomrule
\end{tabular}
\end{center}

\end{document}
